% =====================================================================
%  Rupture de symetrie (lex-leader) -- explication detaillee
% =====================================================================
\documentclass[11pt,a4paper]{article}

\usepackage[utf8]{inputenc}
\usepackage[T1]{fontenc}
\usepackage{lmodern}
\usepackage[french]{babel}

\usepackage[margin=2.2cm]{geometry}
\usepackage{amsmath,amssymb}
\usepackage{booktabs}
\usepackage{enumitem}
\setlist[itemize,1]{label=$\bullet$}
\usepackage{xcolor}
\usepackage{listings}
\usepackage[hidelinks]{hyperref}

\lstset{
  basicstyle=\ttfamily\small,
  frame=single,
  framesep=4pt,
  backgroundcolor=\color{gray!7},
  columns=fullflexible,
  breaklines=true,
}

\hypersetup{
  pdftitle={Rupture de symetrie (lex-leader) -- explication detaillee},
  pdfauthor={T. N. Ramaherisoa},
}

\newcommand{\GD}{\ensuremath{G_D}}

\title{\vspace{-1.5em}Rupture de symétrie (lex-leader) -- de A à Z}
\author{}
\date{}

\begin{document}
\maketitle
\vspace{-2em}

\section{Le problème de départ}

Une molécule benzénoïde donnée (un graphe dual \GD, un sommet par cycle)
peut avoir des \textbf{symétries géométriques} -- rotations, réflexions.
Si on substitue des tailles $5/6/7$ aux sommets sans tenir compte de ça,
le solveur va trouver, pour une substitution donnée, \textbf{toutes ses
copies symétriques} comme des solutions distinctes -- alors que c'est
chimiquement la \textbf{même molécule}. Sans correction, on compterait
(et énumérerait) des doublons inutiles.

\section{Étape 1 -- détecter les automorphismes (bibliothèque)}

Un \textbf{automorphisme} de \GD\ est une permutation $\pi$ des sommets
qui \textbf{préserve toutes les arêtes} : si $i,j$ sont voisins, alors
$\pi(i),\pi(j)$ doivent aussi être voisins. C'est exactement la
traduction graphique d'une symétrie géométrique de la molécule -- pas
besoin de connaître la géométrie, l'automorphisme du graphe suffit.

Concrètement, dans \texttt{preprocessing.py::compute\_automorphisms} :

\begin{lstlisting}
G = graph.dual
matcher = nx.algorithms.isomorphism.GraphMatcher(G, G)
for iso in matcher.isomorphisms_iter():
    ...
\end{lstlisting}

On compare \GD\ à \textbf{lui-même} avec l'algorithme \textbf{VF2}
(implémenté dans NetworkX) : il essaie, sommet par sommet, toutes les
correspondances bijectives qui respectent les arêtes, et les énumère
toutes. On exclut ensuite l'identité (la permutation qui ne bouge
rien), qui n'apporte aucune information de symétrie.

À ce stade, on a une liste potentiellement longue : \textbf{tous} les
automorphismes non triviaux de \GD.

\section{Étape 2 -- réduire à un ensemble générateur minimal}

Poser une contrainte pour \emph{chaque} automorphisme trouvé serait
redondant : si $\pi_1$ et $\pi_2$ sont déjà connus, leur composée
$\pi_1 \circ \pi_2$ est automatiquement un automorphisme aussi, pas
besoin de la traiter séparément. On cherche donc un sous-ensemble
minimal qui \textbf{engendre} tout le reste par composition.

L'algorithme (glouton), implémenté dans
\texttt{preprocessing.py::\_extract\_generators} :

\begin{lstlisting}
generateurs <- [] ; groupe <- { identite }
pour chaque automorphisme a trouve (hors identite) :
    si a n'est pas deja dans groupe :
        generateurs <- generateurs + [a]
        groupe <- cloture(generateurs)   # toutes les compositions possibles
        si groupe == tous les automorphismes trouves : arreter
retourner generateurs
\end{lstlisting}

La « clôture » (\texttt{\_generate\_group}) compose itérativement les
générateurs entre eux jusqu'à ce qu'aucune nouvelle permutation
n'apparaisse (point fixe). On s'arrête dès que cette clôture couvre
tout ce qu'on a trouvé à l'étape 1 -- donc on ne garde que le strict
nécessaire, $\{\pi_1,\dots,\pi_m\}$, avec $m$ généralement petit
(souvent 1 ou 2 pour ces graphes) même si le groupe complet peut être
grand.

\begin{quote}
\small\textit{Nuance honnête : c'est une heuristique simple, pas
l'algorithme le plus rigoureux pour de très gros groupes (la référence
serait Schreier--Sims) -- mais suffisant ici.}
\end{quote}

\section{Étape 3 -- poser la contrainte}

Dans \texttt{model.py} :

\begin{lstlisting}
for gen in generators:
    permuted = [x[gen[i]] for i in range(h)]
    satisfy(LexIncreasing(x, permuted))
\end{lstlisting}

Pour chaque générateur $\pi$, on construit le vecteur « permuté »
$(x_{\pi(0)}, \dots, x_{\pi(h-1)})$ et on impose :
\[
  (x_0, \dots, x_{h-1}) \;\le_{\text{lex}}\; (x_{\pi(0)}, \dots, x_{\pi(h-1)})
\]

\texttt{LexIncreasing} est une \textbf{contrainte globale} standard
(XCSP3/Choco) : sa portée $S(c)$ est le vecteur entier, sa relation
$R(c)$ l'ensemble des tuples qui respectent cet ordre. On ne code pas
l'algorithme qui la fait respecter pendant la recherche -- ça, c'est
interne au solveur -- on se contente de la déclarer.

\textbf{Effet} : parmi toutes les solutions équivalentes par symétrie
(une orbite entière), seule celle qui est lexicographiquement minimale
satisfait \emph{simultanément} toutes ces inégalités. Les autres sont
rejetées par le solveur -- elles ne sont même pas énumérées.

\section{Le résidu imparfait (et son rattrapage)}

En pratique, Choco propage \texttt{LexIncreasing} un peu moins
strictement qu'ACE (l'autre solveur utilisé) : il laisse parfois passer
${\sim}0{,}26\%$ de solutions qui sont en fait dans la même orbite
qu'une autre déjà acceptée. Un post-filtre Python
(\texttt{model.py::\_dedup\_by\_orbit}) rattrape ça après coup : pour
chaque solution renvoyée, il recalcule toute son orbite (en appliquant
les générateurs récursivement) et ne garde que le représentant minimal
-- exactement le même critère que la contrainte, mais vérifié une
seconde fois côté Python pour garantir l'alignement exact entre les
deux solveurs.

\section*{En une phrase}

Détecter (bibliothèque) $\rightarrow$ réduire à un générateur minimal
(code du projet) $\rightarrow$ poser une inégalité lex par générateur
(déclaratif, contrainte globale) $\rightarrow$ laisser le solveur
filtrer pendant la recherche $\rightarrow$ re-vérifier après coup pour
combler l'imperfection de la propagation.

\end{document}
